% Created 2021-03-31 mié 01:34
% Intended LaTeX compiler: pdflatex
\documentclass[letterpaper, 12pt]{article}
\usepackage[hidelinks]{hyperref}
\usepackage[left=2.5cm, top=2.5cm, right=2.5cm, bottom=2.5cm]{geometry}
\setcounter{secnumdepth}{0}
\usepackage[spanish, es-tabla, es-nodecimaldot]{babel}
\usepackage[utf8]{inputenc}
\usepackage{graphicx}
\usepackage{caption}
\usepackage{subcaption}
\usepackage{physics}
\usepackage{amsmath}
\usepackage{gensymb}
\usepackage{siunitx}
\usepackage{listings}
\usepackage{pdfpages}
\usepackage{cite} % referencias y bibliografía
\usepackage{babelbib}
\usepackage{enumerate}
\addto{\captionsspanish}{\renewcommand*{\listfigurename}{Índice de Figuras}}
\addto{\captionsspanish}{\renewcommand*{\contentsname}{Índice}}
\addto{\captionsspanish}{\renewcommand*{\listtablename}{Índice de Tablas}}
\pagestyle{plain}
\usepackage{eso-pic}
\usepackage[hidelinks]{hyperref}
\providecommand\phantomsection{}
\usepackage{cleveref}
\newcommand{\parg}{\par \hfill \break}
\newcommand{\tarea}{Tarea 02: Demostraciones de propiedades de
los algoritmos BFS y DFS}
\newcommand{\fecha}{30-marzo-2021}
\date{\today}
\title{Tarea 02}
\hypersetup{
 pdfauthor={Marco Antonio Cardoso Moreno},
 pdftitle={Tarea 02},
 pdfkeywords={},
 pdfsubject={},
 pdfcreator={Emacs 27.2 (Org mode 9.5)}, 
 pdflang={Spanish}}
\begin{document}

\begin{titlepage}
\begin{center}
  \AddToShipoutPictureBG*{
    \includegraphics[width=\paperwidth,height=\paperheight]
        {../../portada/portada.pdf}}
        \vspace*{2mm}
        \rule{\textwidth}{1.5pt}\\
        \vspace*{5mm}
        \textbf{\Large INSTITUTO POLITÉCNICO NACIONAL}\\
        \vspace*{8mm}
        {\Large Centro de Investigación en Computación}\\
        \vspace{8mm}
        {\Large Maestría en Ciencias en Ingeniería de Cómputo}\\
        \vspace{19mm}
        {\Large Diseño y Análisis de Algoritmos}\\
        \vspace{49mm}




        \textbf{\large \tarea}\\





        \vspace*{48mm}
        {\large Docente: Dr. Rolando Quintero Téllez}\\
        \vspace*{20mm}
        {\large Alumno: Cardoso Moreno Marco Antonio}\\
        \vspace*{10mm}
        \rule{\textwidth}{1.5pt}\\
        {\large Semestre A21} \hfill {\large Fecha: \fecha}
    \end{center}
\end{titlepage}

\section{Demostrar}
\label{sec:orgdec23af}
\subsection{Sea \emph{T} el árbol BFS del grafo \(G=(V, E)\), y sea \((x, y)\) una arista en \emph{T}, entonces el nivel de x y y difiere en 1.}
\label{sec:orgdcd39f4}
Abordando la demostración mediante contradicción, supongamos que:
\begin{itemize}
\item \emph{x} pertenece a la capa \(L_i\), y que \emph{y} pertenece a la capa \(L_j\)
\item \(\abs{i - j} > 1\), es decir, que la distancia entre capas es mayor a 1
\item se llega al punto en el que todas las aristas incidentes a \emph{x} se han examinado
\item los nodos descubiertos desde \emph{x} deben pertenecer a la capa \(L_{i+1}\) o a
alguna capa anterior.
\item por lo tanto, si \emph{y} es un nodo vecino de \emph{x}, debe haber sido descubierto en
este momento como última instancia, o antes.
\end{itemize}
\subsection{Sea \(T\) un árbol DFS, sean \(u\) y \(v\) nodos en \(T\) y sea \((u, v)\) una arista en el grago \(G\) que no existe en \(T\). Entonces \(u\) es ancestro de \(v\) o \(v\) es ancestro de \(u\)}
\label{sec:org81c4307}
\begin{itemize}
\item Si \((u, v)\) es una arista de \(G\) que no está presente en \(T\)
\item Asumamos a su vez que primero se llega al nodo \(u\)
\item Durante una llamada recursiva a la función \(DFS(u)\), se analiza la arista
\((u, v)\), ésta no es agregada al árbo \(T\) porque \(v\) es marcado como
\guillemotleft{}explorado\guillemotright{} duante esta llamada recursiva
\item por lo tanto, si el nodo \(v\) no estaba marcado como \guillemotleft{}explorado\guillemotright{} al iniciar el
llamado recursivo \(DFS(u)\), pero si al final de éste, se infiere que el nodo \(v\) fue descubierto durante \(DFS(u)\).
\item retomando lo explicado en la presentación \guillemotleft{}ADA - Unidad 3\guillemotright{}, una de las
propiedades de la búsqueda en profundidad es:
\begin{itemize}
\item \guillemotleft{}Para una llamada \(DFS(u)\), todos los nodos marcados como \emph{explorados} entre
la invocación y el regreso, son  descendientes de \(u\) en el árbol \(T\)\guillemotright{}
\end{itemize}
\item de los últimos dos puntos, se concluye que \(v\) es un descendiente de \(u\)
\end{itemize}
\end{document}
